% ========================================================
% 2related_work.tex
% ========================================================
\section{Related Work}
\label{sec:related}

\subsection{Spatio-Temporal Load Prediction for Wireless Networks}

Spatio-temporal prediction of network traffic has evolved from standalone LSTM
models~\cite{Shaabanzadeh2024} through graph-enhanced
architectures~\cite{Jin2024} to multi-grained hybrid
approaches~\cite{Liu2024}. Shaabanzadeh \textit{et al.}~\cite{Shaabanzadeh2024}
showed that Transformers outperform CNN and LSTM on campus WiFi traces with
100+ APs, while GraFSTNet~\cite{GraFSTNet2025} introduced dual-branch
spatial-temporal fusion for cellular traffic. However, none of these works
exploit structured event metadata (schedules, timetables) as cross-attention
features, limiting performance during sudden, schedule-driven load transitions.
Salami \textit{et al.}~\cite{Salami2024} explored federated LSTM with
knowledge distillation for per-AP prediction but omit spatial graph context
and schedule awareness. \textbf{STAG-FedSAC's ST-GCAT is the first predictor
to bring schedule features into a graph cross-attention mechanism.}

\subsection{Deep Reinforcement Learning for WiFi Resource Allocation}

DDPG-based control of multi-AP dense WiFi deployments was studied by
Zhang \textit{et al.}~\cite{Zhang2024}, who demonstrated cooperative access
control benefits compared with static association. D3PG~\cite{Liu2024DRL}
replaces DDPG with a diffusion policy for smoother action spaces.
Szott \textit{et al.}~\cite{Szott2024} applied DQN and DDPG to
IEEE~802.11ax contention window optimisation. SAC has emerged as superior
for non-stationary wireless environments due to entropy
regularisation~\cite{SACvsRIS2024}, and was adopted for NB-IoT
fairness-energy trade-offs~\cite{SACNBIoT2024}. These prior works, however,
treat power, channel, and bandwidth as separate or flat action spaces and do
not augment the state with spatio-temporal graph embeddings from a predictor.

\subsection{Prediction--Control Coupling in Wireless Optimisation}

Integrating traffic forecasting with DRL has been studied in O-RAN
slicing~\cite{Lotfi2024} and DQN with lookahead
states~\cite{CNNLSTMDQN2025}. A key finding~\cite{ForecastDRL2024} is
that a frozen predictor used as state augmentation degrades with prediction
errors. STAG-FedSAC addresses this by jointly training predictor and policy,
creating a feedback loop absent in prior work.

\subsection{Federated Deep Reinforcement Learning}

Federated DRL has been applied to contention window
optimisation~\cite{Tu2025}, vehicular
networks~\cite{AFL2025}, and energy-aware IoT
scheduling~\cite{PerFedSAC2025}. E-FRL~\cite{Du2024} is the most relevant
prior work: it applies federated DDPG to dense WiFi channel access with
exponential weighted aggregation. STAG-FedSAC extends E-FRL by adding a
two-level hierarchy, schedule-conditioned prediction, joint training, and
knowledge distillation compression. Tu \textit{et al.}~\cite{Tu2025} proposed
multi-agent DDPG with federated parameter sharing but without prediction-to-control
coupling or personalised local heads. KD-AFRL~\cite{KDAFRL2025} uses KD for
communication efficiency in IoT but without hierarchical zone-level aggregation.

\subsection{Fairness in Multi-AP Wireless Networks}

Jain's Fairness Index (JFI) is a standard metric for evaluating multi-user
equity~\cite{Gurewitz2024}. Primal-dual (Lagrangian) methods have been
applied to enforce capacity constraints in DDPG-based network
slicing~\cite{PDDPG2024}, and fairness-constrained RL has been formalised
through multi-agent MDP frameworks~\cite{FairMDP2024}. STAG-FedSAC adopts
Lagrangian multipliers for capacity constraints while directly incorporating
JFI into its reward function with per-class QoS thresholds.

\subsection{Summary and Positioning}

Table~\ref{tab:related} summarises the most relevant prior works.
STAG-FedSAC is unique in combining all five desiderata---schedule-aware
prediction, graph-embedded DRL state, joint training, hierarchical
federation, and QoS-differentiated fairness---in a single framework.

\begin{table}[htbp]
\centering
\caption{Comparison of STAG-FedSAC with related work.}
\label{tab:related}
\small
\begin{tabular}{lccccc}
\toprule
\textbf{Method} & \textbf{Schedule} & \textbf{Graph} & \textbf{Joint} & \textbf{Hier.\ Fed.} & \textbf{QoS Fair.} \\
\midrule
Shaabanzadeh~\cite{Shaabanzadeh2024} & \checkmark & $\times$ & $\times$ & $\times$ & $\times$ \\
Zhang (DDPG)~\cite{Zhang2024}        & $\times$ & $\times$ & $\times$ & $\times$ & $\times$ \\
E-FRL~\cite{Du2024}                  & $\times$ & $\times$ & $\times$ & $\times$ & $\times$ \\
Salami (FedLSTM)~\cite{Salami2024}   & $\times$ & $\times$ & $\times$ & $\times$ & $\times$ \\
Lotfi (O-RAN)~\cite{Lotfi2024}       & $\times$ & $\times$ & Partial & $\times$ & $\times$ \\
PerFedSAC~\cite{PerFedSAC2025}       & $\times$ & $\times$ & $\times$ & $\times$ & \checkmark \\
\textbf{STAG-FedSAC (ours)}          & \checkmark & \checkmark & \checkmark & \checkmark & \checkmark \\
\bottomrule
\end{tabular}
\end{table}
